% !TEX root = ../../main.tex
% C1.2 — Mục tiêu của đề tài
% Chương 1 là tổng quan: nêu mục tiêu ở mức "đạt được gì", KHÔNG liệt kê chức năng chi tiết, trường dữ liệu hay bộ lọc
% (để dành cho Mục 4.2), KHÔNG nêu tên công nghệ/mô hình (Chương 3, Mục 6.1) và KHÔNG nêu tên vai trò (Mục 4.1).
% Mỗi mục tiêu phải được đối chiếu lại ở Mục 9.1.
\section{Mục tiêu của đề tài}
\label{sec:1.2}

Mục tiêu tổng quát là xây dựng một ứng dụng hỗ trợ tìm kiếm người trong dữ liệu camera giám sát bằng nhiều hình thức mô tả thay vì xem lại video thủ công. Hệ thống đề xuất và xếp hạng kết quả; quyết định một kết quả có đúng là người cần tìm hay không thuộc về người dùng.

Các mục tiêu cụ thể gồm:

\begin{enumerate}
    \item \textbf{Chuyển dữ liệu camera thành dữ liệu có thể tìm kiếm.} Tự động phân tích hình ảnh camera, phát hiện và theo vết người, và lưu mỗi lần xuất hiện như một đơn vị tra cứu được.

    \item \textbf{Tìm kiếm người bằng mô tả đa phương thức.} Mô tả người cần tìm bằng ảnh tham chiếu, câu mô tả hoặc đặc điểm ngoại hình chọn sẵn, cùng được so khớp với dữ liệu đã phân tích và xếp theo mức độ phù hợp.

    \item \textbf{Hỗ trợ đánh giá và quản lý kết quả tìm kiếm.} Xem từng kết quả trong bối cảnh gốc để tự đánh giá, và lưu kết quả đã xác nhận thành hồ sơ vụ việc để theo dõi.

    \item \textbf{Vận hành an toàn trong môi trường nhiều người dùng.} Phân quyền theo trách nhiệm, giới hạn phạm vi tìm kiếm của mỗi người theo khu vực được giao, và cung cấp chức năng quản trị, theo dõi hệ thống.

    \item \textbf{Đánh giá khả năng áp dụng các mô hình trí tuệ nhân tạo có sẵn.} Không huấn luyện mô hình mới mà chọn các mô hình đã công bố phù hợp với bài toán và phần cứng phổ thông, rồi đánh giá chất lượng tìm kiếm và chi phí xử lý khi tích hợp vào ứng dụng.
\end{enumerate}

Đề tài được xem là hoàn thành mục tiêu khi ứng dụng vận hành trọn vẹn luồng nghiệp vụ, từ tiếp nhận dữ liệu camera, tìm kiếm, đánh giá kết quả đến lưu hồ sơ vụ việc, đồng thời có số liệu đánh giá chất lượng tìm kiếm và hiệu năng được trình bày tại Chương~\ref{chapter:kiemthu}.
